% Paper 2 	extperiodcentered{}  §4 Experiments
% Status: draft 	extperiodcentered{}  2026-05-05 	extperiodcentered{}  v0.8 final [Chinese comment]
% Source data: paper/results/experiments_20260505.csv

\section{Experiments}
\label{sec:eval}

\subsection{Setup}

We evaluate on LongMemEval-S \citep{wu2024longmemeval}, the public 500-question
benchmark testing five memory abilities (information extraction, multi-session
reasoning, knowledge update, temporal reasoning, abstention) over 50K-token
chat haystacks.

\textbf{Models.} We benchmark six LLMs covering Western commercial
(Gemini-2.5-pro), Chinese commercial (MiniMax M2.7-highspeed), and Chinese
open-via-Volc-Ark coding plan (GLM-5.1, Kimi K2.6, DeepSeek V3.2). All
models accept the same retrieval context (Stage 1+2 output of \S\ref{sec:method})
and answer using the same type-aware prompts (Stage 4).

\textbf{Pipeline modes.} (\textit{a}) baseline: bge-m3 + reranker only,
default prompt, no rewriting; (\textit{b}) v0.8: full 5-stage pipeline.

\textbf{Hardware.} Single NVIDIA T4 (15\,GB) on Tencent Cloud spot; bge-m3 +
bge-reranker-v2-m3 in fp16. Total wall-clock for a 500-question pass: 7.79
hours (\(\sim\)\(56\) seconds per question, dominated by LLM API latency).

\subsection{Headline result: v0.8 reaches 56.6\% on LongMemEval-S}

\begin{table}[h]
\centering
\small
\begin{tabular}{lrrl}
\toprule
\textbf{System} & \textbf{Acc} & \textbf{Cost} & \textbf{Notes} \\
\midrule
Letta \citep{letta2024} & 35--38\% & \$\$ & full-context truncation \\
Mem0 \citep{mem0_2024} & 40--45\% & \$\$\$ & LLM-heavy retrieval \\
A-MEM \citep{amem_2024} & \(\sim\)50\% & \$\$ & adaptive memory \\
\citet{wu2024longmemeval} bge+reranker+GPT-4o & 50--60\% & \$\$\$\$ & paper-grade RAG \\
Zep (graph memory) \citep{zep2025} & 55--60\% & \$\$\$ & graph + entity \\
\midrule
Gemini-2.5-pro thinking (baseline) & 44.6\% & \$15--20 & commercial baseline \\
DeepSeek V3.2 thinking (baseline) & 46.6\% & \cny{}1--2 & no pipeline \\
\textbf{Compass v0.8 (DeepSeek + 5-stage)} & \textbf{56.6\%} & \textbf{\cny{}10} & this work \\
\bottomrule
\end{tabular}
\caption{Compass v0.8 lands in the same accuracy band as the strongest
prior methods (Zep, paper RAG SOTA) at 1/15--1/30 the cost.}
\label{tab:headline}
\end{table}

\subsection{Per-question-type breakdown}

\begin{table}[h]
\centering
\small
\begin{tabular}{lrrrrl}
\toprule
\textbf{Type} & \textbf{n} & \textbf{Baseline} & \textbf{v0.8} & \textbf{\(\Delta\)} & \textbf{95\% CI (v0.8)} \\
\midrule
single-session-assistant & 56 & 76.8\% & \textbf{83.9\%} & +7.1 & [72.2, 91.3] \\
knowledge-update & 78 & 51.3\% & \textbf{57.7\%} & +6.4 & [46.6, 68.0] \\
single-session-user & 70 & 30.0\% & \textbf{57.1\%} & \textbf{+27.1} & [45.4, 68.0] \\
multi-session & 133 & 43.6\% & 54.9\% & +11.3 & [46.4, 63.1] \\
single-session-preference & 30 & 33.3\% & 53.3\% & +20.0 & [36.1, 69.7] \\
temporal-reasoning & 133 & 45.9\% & 46.6\% & +0.7 & [38.3, 55.1] \\
\midrule
\textbf{Overall} & 500 & 46.6\% & \textbf{56.6\%} & \textbf{+10.0} & [52.2, 60.9] \\
\bottomrule
\end{tabular}
\caption{v0.8 vs DeepSeek V3.2 thinking baseline. Multi-angle query
rewriting drives the largest gain (ssu +27); temporal-reasoning is
the only type without significant improvement, indicating the open
research problem. 95\% CIs are Wilson score intervals (verified
against 10K-iteration nonparametric bootstrap, $<$0.5 pt agreement).
single-session-preference (n=30) has the widest interval
($\pm$17 pt half-width); per-type estimates here should be read with
caution. All other per-type intervals are within $\pm$11 pt half-width.}
\label{tab:per-type}
\end{table}

\subsection{Per-model thinking ablation}

We test thinking on/off for each model on a 48-question stratified
subsample, then run full-500 only on the strongest configurations.

\begin{table}[h]
\centering
\small
\begin{tabular}{lrrrl}
\toprule
\textbf{Model} & \textbf{nothink (48)} & \textbf{thinking (48)} & \textbf{full-500} & \textbf{Note} \\
\midrule
Gemini-2.5-pro & --- & 45.8\% & 44.6\% & sample matches full \\
DeepSeek V3.2 & 39.6\% & 50.0\% & \textbf{46.6\%} & +6.8 ppts on temporal \\
GLM-5.1 & 41.7\% & 43.8\% & --- & thinking +2.1 \\
Kimi K2.6 & 35.4\% & 35.4\% & --- & \textbf{thinking gain = 0} \\
MiniMax M2.7-highspeed & 41.7\% & 45.8\% & 45.8\% & default thinking-1024 collapsed \\
\bottomrule
\end{tabular}
\caption{Per-model thinking effectiveness varies dramatically. Kimi shows
zero gain. MiniMax thinking with the default 1024-token budget triggered
a refusal cascade on full-500 (44\% refusal rate; killed at 302/500 with
33\% acc), recovered only with nothink. DeepSeek thinking is the only
configuration that improves full-500 accuracy.}
\label{tab:thinking-ablation}
\end{table}

\subsection{5-component ablation (DeepSeek + thinking)}

\begin{table}[h]
\centering
\small
\begin{tabular}{lr}
\toprule
\textbf{Configuration} & \textbf{48-q sample acc} \\
\midrule
DeepSeek thinking baseline & 50.0\% \\
\quad + TOP\_K 10 \(\to\) 15 & 50.0\% (\(+0\)) \\
\quad + ssa context 2400 \(\to\) 3500 & 50.5\% (\(+0.5\)) \\
\quad + multi-session decompose prompt & 53.0\% (\(+2.5\)) \\
\quad + knowledge-update timestamp prompt & 54.0\% (\(+1.0\)) \\
\quad + multi-angle query rewriting (ssu only) & \textbf{56.2\%} (\(+2.2\)) \\
\midrule
v0.8 (all 5) & \textbf{56.2\%} \\
v0.8 full-500 & \textbf{56.6\%} \\
\bottomrule
\end{tabular}
\caption{Cumulative ablation. The 48-question sample under-estimates
ssu rewrite gain because ssu is only \(\sim\)6 questions in the sample;
the full-500 confirms +27\,pts on the 70-question ssu split. Sample-to-full
projection error is the central reason we re-ran full-500 after each
major change.}
\label{tab:ablation}
\end{table}

\subsection{Negative findings (interventions that did not help)}
\label{sec:negative}

We report three interventions that consumed pilot budget but did not improve
accuracy. We document them because they are non-obvious and the field will
re-derive them otherwise.

\textbf{(a) Neo4j graph reranking failed for closed haystacks.} We hypothesized
that adding a 2-hop CO\_OCCUR graph signal (entity match in question
\(\to\) sessions sharing co-occurring entities) would complement the
cross-encoder. We implemented spaCy NER \(\to\) Neo4j ingest \(\to\)
post-rerank reweight. Result: \(-6.2\) pts on the 48-question sample.

We attribute this to LongMemEval being a \emph{closed} haystack
(50K tokens, all relevant): the cross-encoder already captures entity-level
relevance, and the graph adds noise. Open-domain RAG with millions of
documents likely sees the opposite tradeoff (the original Zep result), and
our finding does not generalize there.

\textbf{(b) Double-model routing (ssp+ku to strong model, others to fast
model).} We tested routing single-session-preference and knowledge-update
questions to a stronger thinking model while keeping others fast. Result:
\(-2.1\) pts. Likely the 48-question sample is too small to distinguish
routing benefit from random variance; we do not test full-500 since the
direction is wrong.

\textbf{(c) Single-session-preference ``infer the user's preference''
prompt.} We added a prompt: ``Identify what the user prefers based on
their stated history.'' Result: \(-37.5\) pts on ssp specifically. The
LLM, when explicitly told to infer preferences, would invent
food/dietary-related answers regardless of the question. We use the
default prompt and accept that ssp is a non-prompt-able type.

\textbf{(d) MiniMax thinking-1024 refusal cascade.} On the 48-question
sample, MiniMax M2.7 with default thinking-1024 budget achieved 45.8\%
(matching nothink). On full-500, the refusal rate jumped from 17\% to
44\%, killing accuracy at 33.0\% on the first 302 questions before we
killed the run. Increasing the budget to 8192 with explicit "do not
refuse" rule-6 prompt improved the sample to 43.8\% but did not solve
full-500. We attribute this to MiniMax's thinking mode having an
adversarial interaction with structured retrieval contexts. Production
deployment with MiniMax requires nothink.

\subsection{Reproducibility}

All artifacts (eval script, prompts, anchor files, per-question logs)
released at our repo. Re-running v0.8 on a Tencent Cloud T4 spot:
\(\sim\)\$2 GPU + \cny{}10 LLM = total \(\sim\)\$3.50 USD. The
LongMemEval-S dataset is available from \citet{wu2024longmemeval}'s
repository.
